We compare the two-vector and one-vector ERE backends as the exact layer
of two ARE wrappers (SODA and Seg-ARE) across six distributions.
The analytic prediction from
Section~\ref{subsec:ere-one-vector} is $\Delta = M/n$~BPK, where $M$ is
the number of non-empty blocks; the maximum is $B/n = 1$~BPK and the
Poisson middle for uniform input at $\lambda = 1$ is
$1 - e^{-1} \approx 0.632$~BPK. The numbers below confirm the prediction
end-to-end.

\begin{table}[H]
    \centering
    \setlength{\tabcolsep}{5pt}
    \small
    \begin{tabular}{@{}lrrrrrr@{}}
        \toprule
        & \multicolumn{3}{c}{SODA} & \multicolumn{3}{c}{Seg-ARE} \\
        \cmidrule(lr){2-4} \cmidrule(lr){5-7}
        Dataset
        & two-vec BPK & one-vec BPK & $\Delta$
        & two-vec BPK & one-vec BPK & $\Delta$ \\
        \midrule
        \texttt{uniform}     & $16.932$ & $16.499$ & $0.433$ & $16.932$ & $16.499$ & $0.433$ \\
        \texttt{clustered}   & $16.602$ & $16.500$ & $0.102$ & $12.759$ & $12.519$ & $0.240$ \\
        \texttt{sosd\_fb}    & $16.020$ & $16.000$ & $0.020$ & $11.232$ & $11.000$ & $0.232$ \\
        \texttt{sosd\_wiki}  & $16.532$ & $16.530$ & $0.002$ & $9.708$  & $9.530$  & $0.178$ \\
        \texttt{sosd\_osm}   & $16.931$ & $16.499$ & $0.432$ & $26.678$ & $26.411$ & $0.267$ \\
        \texttt{sosd\_books} & $16.000$ & $16.000$ & $0.000$ & $4.517$  & $4.000$  & $0.517$ \\
        \bottomrule
    \end{tabular}
    \caption{Total ARE-level BPK with the two-vector $(D_1, D_2)$ ERE backend
        (Section~\ref{subsec:ere-original-structure}) versus the one-vector $D$
        backend (Section~\ref{subsec:ere-one-vector}). Workload: $n = 2^{20}$
        (or all available keys, whichever is smaller), $K = 34$ — equivalent to
        $\mathcal{L} = 128$, $\varepsilon = 0.01$ via $K = \lceil \log_2(n\mathcal{L}/\varepsilon) \rceil$
        for SODA; Seg-ARE takes $K$ directly. SODA stays in the
        $14$--$16$~BPK headline window (Table~\ref{tab:industry-bpk}) on every
        distribution; Seg-ARE ranges from $4.0$~BPK (\texttt{sosd\_books},
        single-cluster collapse) to $26.4$~BPK (\texttt{sosd\_osm}, many small
        clusters with per-cluster bookkeeping overhead).}
    \label{tab:ere-backend-cross-distribution}
\end{table}

SODA's locality-preserving hash keeps clustered inputs clustered in the
hashed universe, so the SOSD datasets leave most blocks empty under
SODA ($\Delta \le 0.020$~BPK on \texttt{sosd\_fb}, \texttt{sosd\_wiki},
\texttt{sosd\_books}); \texttt{uniform} and \texttt{sosd\_osm} approach the analytic prediction
(Section~\ref{subsec:block-sizes-and-lookup-strategy}) at ${\sim}0.43$~BPK
(because the benchmark evaluates at exactly $n = 2^{20}$,
SODA hash collisions drop the unique fingerprint count slightly below the power-of-two threshold;
this halves the ERE block count $B$ to $2^{19}$, shifting the load factor to $\lambda \approx 2$;
the reported ARE BPK is logical, so it perfectly matches the analytic $M/n$ prediction without physical succinct-index overheads).
Seg-ARE isolates dense regions and
builds a separate ERE per cluster; even within a cluster the keys
retain sub-structure, so $M/B$ stays well below $1$. Its absolute BPK
is dominated by the wrapper rather than ERE, so even small $M/n$
values translate into a noticeable relative improvement (e.g.\
\texttt{sosd\_books} drops from $4.517$ to $4.000$~BPK).

\paragraph{Query latency.}
Table~\ref{tab:ere-backend-query-latency} reports average ARE query
latency on the same workload. Both wrappers gain uniformly across all
distributions: SODA $1.05$--$3.59\times$, Seg-ARE $1.28$--$3.00\times$.
The $\geq 2\times$ wins on \texttt{uniform} and \texttt{sosd\_osm} are
instruction-driven: SODA's
2-universal hashing maps both to a near-uniform image in $[0, 2^K)$,
so the merged $D$ vector's \textsc{Select} chains are short. The smaller speedups on
\texttt{clustered}, \texttt{sosd\_fb}, and \texttt{sosd\_wiki} match
their $34$--$39\%$ instruction reduction; \texttt{sosd\_books} gains
least because its power-law distribution leaves many hashed buckets
empty even after hashing, lengthening \textsc{Select} chains and
limiting the instruction saving to $12\%$. At $n = 2^{20}$, cache miss
counts stay below one per query---the filter fits in the L3 cache---so
cache effects are secondary.

\begin{table}[H]
    \centering
    \setlength{\tabcolsep}{6pt}
    \small
    \begin{tabular}{@{}lrrr@{}}
        \toprule
        Dataset & $\Delta$~instrs & $\Delta$~L1-loads & $\Delta$~L1-misses \\
        \midrule
        \texttt{uniform}     & $-76\%$ & $-81\%$ & $-56\%$ \\
        \texttt{clustered}   & $-34\%$ & $-35\%$ & $-20\%$ \\
        \texttt{sosd\_fb}    & $-39\%$ & $-39\%$ & $-21\%$ \\
        \texttt{sosd\_wiki}  & $-39\%$ & $-38\%$ & $-14\%$ \\
        \texttt{sosd\_osm}   & $-76\%$ & $-81\%$ & $-56\%$ \\
        \texttt{sosd\_books} & $-12\%$ & $-13\%$ & $\phantom{-}4\%$  \\
        \bottomrule
    \end{tabular}
    \caption{Per-query reduction in instructions executed, L1 cache loads,
        and L1 cache misses when switching from the two-vector to the
        one-vector ERE backend inside SODA ARE. Measured with
        \texttt{perf\_event\_open} (user-space only) on an x86-64 Linux
        machine (Linux~6.17, 62\,GB RAM), pinned to a single core;
        $n = 2^{20}$, $K = 34$, range length $128$,
        100\,000 queries (10\,000 warmup).}
    \label{tab:ere-backend-pmu}
\end{table}

\begin{table}[H]
    \centering
    \setlength{\tabcolsep}{5pt}
    \small
    \begin{tabular}{@{}lrrrrrr@{}}
        \toprule
        & \multicolumn{3}{c}{SODA} & \multicolumn{3}{c}{Seg-ARE} \\
        \cmidrule(lr){2-4} \cmidrule(lr){5-7}
        Dataset
        & two-vec ns & one-vec ns & speedup
        & two-vec ns & one-vec ns & speedup \\
        \midrule
        \texttt{uniform}     & $324.1$ & $90.6$  & $3.58\times$ & $218.6$ & $72.9$  & $3.00\times$ \\
        \texttt{clustered}   & $381.8$ & $287.5$ & $1.33\times$ & $317.9$ & $198.3$ & $1.60\times$ \\
        \texttt{sosd\_fb}    & $445.6$ & $309.5$ & $1.44\times$ & $224.9$ & $158.0$ & $1.42\times$ \\
        \texttt{sosd\_wiki}  & $424.7$ & $327.3$ & $1.30\times$ & $303.3$ & $206.5$ & $1.47\times$ \\
        \texttt{sosd\_osm}   & $334.9$ & $93.4$  & $3.59\times$ & $298.3$ & $232.6$ & $1.28\times$ \\
        \texttt{sosd\_books} & $323.7$ & $308.4$ & $1.05\times$ & $223.1$ & $104.7$ & $2.13\times$ \\
        \bottomrule
    \end{tabular}
    \caption{Average ARE query latency, two-vector vs one-vector ERE backend.
    Same workload as Table~\ref{tab:ere-backend-cross-distribution}:
        $n = 2^{20}$ (or all available keys, whichever is smaller), $K = 34$;
        raw $64$-bit keys (no $60$-bit mask). x86-64 Linux machine
        (Linux~6.17, 62\,GB RAM), pinned to a single core;
        100\,000 queries (10\,000 warmup).
        This is a small-scale backend
        comparison at fixed $n = 2^{20}$; absolute latencies grow $5$--$10\times$
        at the headline scale $n = 2^{24}$ once the filter footprint outgrows
        the L3 cache and per-query work becomes cache-miss-bound.}
    \label{tab:ere-backend-query-latency}
\end{table}